\section{Interface screenshots of every view}\label{app:shots}
All interface screenshots of the project are collected here, grouped as the result views of Section~\ref{sec:views}. Each is a full-window capture of SWAFL Desktop (version 1.1.0, window of 2000$\times$1125 points at twice the pixel density); the full-resolution files (4000$\times$2250\,px) are in \texttt{Manuscripts/figures/shots/} of the repository, and every screen is documented, with numbered arrows on each element, in the \manuallink.
\begin{figure}[H]\centering
\begin{subfigure}{\textwidth}\centering\includegraphics[width=0.86\textwidth,height=0.4\textheight,keepaspectratio]{\figdir/shots_jpg/layer_field.jpg}\caption{Field balance, the home view: daily fluxes on the field clock, the soil profile with sensor, roots and wetting front, and the closed budget (\cref{eq:balance}).}\end{subfigure}\\[4pt]
\begin{subfigure}{\textwidth}\centering\includegraphics[width=0.86\textwidth,height=0.4\textheight,keepaspectratio]{\figdir/shots_jpg/layer_spac.jpg}\caption{Soil, plant and air: layouts of the five irrigation methods and the water potential along the path from the deep soil to the air (\cref{eq:psisoil,eq:spac,eq:psiair}).}\end{subfigure}
\caption{SWAFL Desktop: field and the soil--plant--air path.}\label{fig:shot-layer-field}
\end{figure}
\begin{figure}[H]\centering
\begin{subfigure}{\textwidth}\centering\includegraphics[width=0.86\textwidth,height=0.4\textheight,keepaspectratio]{\figdir/shots_jpg/layer_v1.jpg}\caption{Sensor and trigger: 20 years of daily suction at the sensor, irrigated and rainfed, against the trigger, and the days beyond the trigger per year (\cref{eq:Ny}).}\end{subfigure}\\[4pt]
\begin{subfigure}{\textwidth}\centering\includegraphics[width=0.86\textwidth,height=0.4\textheight,keepaspectratio]{\figdir/shots_jpg/layer_v2.jpg}\caption{Profile evolution: pressure-head profile through the season, irrigated and rainfed, and the share of the root zone beyond the trigger (\cref{eq:share}).}\end{subfigure}
\caption{SWAFL Desktop: soil water and irrigation timing.}\label{fig:shot-layer-v1}
\end{figure}
\begin{figure}[H]\centering
\begin{subfigure}{\textwidth}\centering\includegraphics[width=0.86\textwidth,height=0.4\textheight,keepaspectratio]{\figdir/shots_jpg/layer_n1.jpg}\caption{Trigger law: root-zone depletion of ten soils on one axis with the critical depletion $f^\ast$, each soil's own trigger suction and yield, and the gauge of the sensor through the season (\cref{eq:depl,eq:invdepl,eq:Zn})}\end{subfigure}
\caption{SWAFL Desktop: trigger law.}\label{fig:shot-layer-n1}
\end{figure}
\begin{figure}[H]\centering
\begin{subfigure}{\textwidth}\centering\includegraphics[width=0.86\textwidth,height=0.4\textheight,keepaspectratio]{\figdir/shots_jpg/layer_v3.jpg}\caption{Tillage and structure: modelled thin section of the topsoil with the water- and air-filled pores at the chosen suction, the pore sizes, and the same place under conventional, reduced and no tillage (\cref{eq:jurin})}\end{subfigure}\\[4pt]
\begin{subfigure}{\textwidth}\centering\includegraphics[width=0.86\textwidth,height=0.4\textheight,keepaspectratio]{\figdir/shots_jpg/layer_v7.jpg}\caption{Soil hydraulic drivers: tillage and consolidation effects on drainage for the ten soils and their regressions on clay content (\cref{eq:drivers}).}\end{subfigure}
\caption{SWAFL Desktop: soil structure and hydraulic drivers.}\label{fig:shot-layer-v3}
\end{figure}
\begin{figure}[H]\centering
\begin{subfigure}{\textwidth}\centering\includegraphics[width=0.86\textwidth,height=0.4\textheight,keepaspectratio]{\figdir/shots_jpg/layer_r1.jpg}\caption{Root structure: rooting depth and the depths holding 50 and 90\,\% of uptake in the ten soils (\cref{eq:rmax,eq:rd,eq:rootdens,eq:zp}), with the E15 rooting-depth heat maps.}\end{subfigure}
\caption{SWAFL Desktop: roots.}\label{fig:shot-layer-r1}
\end{figure}
\begin{figure}[H]\centering
\begin{subfigure}{\textwidth}\centering\includegraphics[width=0.86\textwidth,height=0.4\textheight,keepaspectratio]{\figdir/shots_jpg/layer_e1.jpg}\caption{Evaporation from the pore to the atmosphere: a modelled section of coarse loam drying after a wetting, the drying front, stage-1 and stage-2 evaporation and the water lost since wetting (\cref{eq:twostage,eq:front})}\end{subfigure}\\[4pt]
\begin{subfigure}{\textwidth}\centering\includegraphics[width=0.86\textwidth,height=0.4\textheight,keepaspectratio]{\figdir/shots_jpg/layer_v4.jpg}\caption{Mulch against bare soil: six mulch materials, and soil evaporation, transpiration, irrigation, drainage and nitrate leaching under bare soil, the chosen mulch and the strongest mulch (\cref{eq:mulchdelta})}\end{subfigure}
\caption{SWAFL Desktop: evaporation and mulch.}\label{fig:shot-layer-e1}
\end{figure}
\begin{figure}[H]\centering
\begin{subfigure}{\textwidth}\centering\includegraphics[width=0.86\textwidth,height=0.4\textheight,keepaspectratio]{\figdir/shots_jpg/layer_v6.jpg}\caption{Mulch control: soil evaporation by mulch factor and interception capacity, and drainage by tillage and consolidation rate.}\end{subfigure}\\[4pt]
\begin{subfigure}{\textwidth}\centering\includegraphics[width=0.86\textwidth,height=0.4\textheight,keepaspectratio]{\figdir/shots_jpg/layer_n2.jpg}\caption{Fate of saved water: evaporation saved, irrigation saved, extra drainage and extra transpiration against the rainfall scale (\cref{eq:fate}).}\end{subfigure}
\caption{SWAFL Desktop: mulch control and the fate of saved water.}\label{fig:shot-layer-v6}
\end{figure}
\begin{figure}[H]\centering
\begin{subfigure}{\textwidth}\centering\includegraphics[width=0.86\textwidth,height=0.4\textheight,keepaspectratio]{\figdir/shots_jpg/layer_n3.jpg}\caption{Regime map: mulch factor against rainfall scale coloured by the dominant loss (\cref{eq:regime}).}\end{subfigure}\\[4pt]
\begin{subfigure}{\textwidth}\centering\includegraphics[width=0.86\textwidth,height=0.4\textheight,keepaspectratio]{\figdir/shots_jpg/layer_v5.jpg}\caption{Feedback loops: the seven loops with arrow widths scaled by the effect of one control step (\cref{tab:loops}).}\end{subfigure}
\caption{SWAFL Desktop: regime map and feedback loops.}\label{fig:shot-layer-n3}
\end{figure}
\begin{figure}[H]\centering
\begin{subfigure}{\textwidth}\centering\includegraphics[width=0.86\textwidth,height=0.4\textheight,keepaspectratio]{\figdir/shots_jpg/layer_sid.jpg}\caption{Surface-irrigation layer: start from your soil or a textbook design (A7, A9, A11, D1), border strip, furrow or basin, intake from the soil or by $k$ and $A$, the cutoff rule, and the inputs carried from the soil, root zone and trigger of the DSS (Section~\ref{sec:surf})}\end{subfigure}\\[4pt]
\begin{subfigure}{\textwidth}\centering\includegraphics[width=0.86\textwidth,height=0.4\textheight,keepaspectratio]{\figdir/shots_jpg/sid_sec_long.jpg}\caption{Surface irrigation: layout in plan with the advance front shaded by the depth already infiltrated, and the longitudinal section with the wetting front, the required depth $Z_n$, deep percolation and runoff (\cref{eq:vb,eq:vols})}\end{subfigure}
\caption{SWAFL Desktop: surface irrigation.}\label{fig:shot-layer-sid}
\end{figure}
\begin{figure}[H]\centering
\begin{subfigure}{\textwidth}\centering\includegraphics[width=0.86\textwidth,height=0.4\textheight,keepaspectratio]{\figdir/shots_jpg/sid_sec_sched.jpg}\caption{Surface irrigation: hydrograph of inflow, infiltration and runoff through time.}\end{subfigure}\\[4pt]
\begin{subfigure}{\textwidth}\centering\includegraphics[width=0.86\textwidth,height=0.4\textheight,keepaspectratio]{\figdir/shots_jpg/sid_sec_table.jpg}\caption{Surface irrigation: design table of options with application and storage efficiency, deep percolation and runoff (\cref{eq:eff}).}\end{subfigure}
\caption{SWAFL Desktop: surface irrigation, schedule and design table.}\label{fig:shot-sid-sec-sched}
\end{figure}
\begin{figure}[H]\centering
\begin{subfigure}{\textwidth}\centering\includegraphics[width=0.86\textwidth,height=0.4\textheight,keepaspectratio]{\figdir/shots_jpg/layer_drip.jpg}\caption{Drip-irrigation layer: crop, mulch, emitter type and placement, the design inputs, and the soil moisture around the driest, the average and the wettest emitter (Section~\ref{sec:drip-m})}\end{subfigure}\\[4pt]
\begin{subfigure}{\textwidth}\centering\includegraphics[width=0.86\textwidth,height=0.4\textheight,keepaspectratio]{\figdir/shots_jpg/drip_sec_layout.jpg}\caption{Drip irrigation: field layout and infrastructure.}\end{subfigure}
\caption{SWAFL Desktop: drip irrigation.}\label{fig:shot-layer-drip}
\end{figure}
\begin{figure}[H]\centering
\begin{subfigure}{\textwidth}\centering\includegraphics[width=0.86\textwidth,height=0.4\textheight,keepaspectratio]{\figdir/shots_jpg/drip_sec_lateral.jpg}\caption{Drip irrigation: head and emitter flow along the lateral (\cref{eq:HW}) with the 20\,\% loss warning.}\end{subfigure}\\[4pt]
\begin{subfigure}{\textwidth}\centering\includegraphics[width=0.86\textwidth,height=0.4\textheight,keepaspectratio]{\figdir/shots_jpg/drip_sec_result.jpg}\caption{Drip irrigation: result summary with emission uniformity (\cref{eq:Eu}) and efficiencies.}\end{subfigure}
\caption{SWAFL Desktop: drip irrigation, lateral and results.}\label{fig:shot-drip-sec-lateral}
\end{figure}
\begin{figure}[H]\centering
\begin{subfigure}{\textwidth}\centering\includegraphics[width=0.86\textwidth,height=0.4\textheight,keepaspectratio]{\figdir/shots_jpg/drip_sec_table.jpg}\caption{Drip irrigation: design table, iterations toward an application efficiency of 85\,\% or more.}\end{subfigure}
\caption{SWAFL Desktop: drip irrigation, design table.}\label{fig:shot-drip-sec-table}
\end{figure}
\begin{figure}[H]\centering
\begin{subfigure}{\textwidth}\centering\includegraphics[width=0.86\textwidth,height=0.4\textheight,keepaspectratio]{\figdir/shots_jpg/layer_spr.jpg}\caption{Sprinkler layer: nozzle, wetted radius, pattern, overlap and uniformity (\cref{eq:nozzle}; Section~\ref{sec:sprk}).}\end{subfigure}\\[4pt]
\begin{subfigure}{\textwidth}\centering\includegraphics[width=0.86\textwidth,height=0.4\textheight,keepaspectratio]{\figdir/shots_jpg/spr_sec_checks.jpg}\caption{Sprinkler design: textbook decision checklist (sprinkler class, spacing, application rate, sets per day).}\end{subfigure}
\caption{SWAFL Desktop: sprinkler irrigation.}\label{fig:shot-layer-spr}
\end{figure}
\begin{figure}[H]\centering
\begin{subfigure}{\textwidth}\centering\includegraphics[width=0.86\textwidth,height=0.4\textheight,keepaspectratio]{\figdir/shots_jpg/spr_sec_intake.jpg}\caption{Sprinkler: layout in plan, with the depth infiltrated from the overlapping spray patterns.}\end{subfigure}\\[4pt]
\begin{subfigure}{\textwidth}\centering\includegraphics[width=0.86\textwidth,height=0.4\textheight,keepaspectratio]{\figdir/shots_jpg/dash_design.jpg}\caption{Dashboard with the design tab open beside the field-balance layer.}\end{subfigure}
\caption{SWAFL Desktop: sprinkler intake and the dashboard.}\label{fig:shot-spr-sec-intake}
\end{figure}
\begin{figure}[H]\centering
\begin{subfigure}{\textwidth}\centering\includegraphics[width=0.86\textwidth,height=0.4\textheight,keepaspectratio]{\figdir/shots_jpg/dash_results.jpg}\caption{Dashboard, results tab: annual and daily flux tiles of the field water balance.}\end{subfigure}\\[4pt]
\begin{subfigure}{\textwidth}\centering\includegraphics[width=0.86\textwidth,height=0.4\textheight,keepaspectratio]{\figdir/shots_jpg/dash_charts.jpg}\caption{Dashboard, charts tab: seasonal rainfall, irrigation, drainage, evaporation and transpiration.}\end{subfigure}
\caption{SWAFL Desktop: dashboard.}\label{fig:shot-dash-results}
\end{figure}
\begin{figure}[H]\centering
\begin{subfigure}{\textwidth}\centering\includegraphics[width=0.86\textwidth,height=0.4\textheight,keepaspectratio]{\figdir/shots_jpg/dash_messages.jpg}\caption{Dashboard, messages tab: the decision messages and warnings of the current state.}\end{subfigure}
\caption{SWAFL Desktop: dashboard messages.}\label{fig:shot-dash-messages}
\end{figure}

\begin{figure}[H]\centering
\begin{subfigure}{0.48\textwidth}\centering\includegraphics[width=\textwidth,height=0.42\textheight,keepaspectratio]{\figdir/shots_jpg/real_field_balance_charts.jpg}\caption{Field balance, charts: evaporation, transpiration and evapotranspiration of 2026 against the median of 20 seasons, with rain, irrigation and canopy cover, and the root-zone water against the trigger with the 10 to 90\,\% range and the irrigation events.}\end{subfigure}\hfill
\begin{subfigure}{0.48\textwidth}\centering\includegraphics[width=\textwidth,height=0.42\textheight,keepaspectratio]{\figdir/shots_jpg/real_crop_water.jpg}\caption{Crop water requirement: crop coefficient and daily need through the year, and the root-zone depletion with the trigger, the readily and the totally available water and the irrigations of 2026.}\end{subfigure}
\caption{SWAFL Desktop: the field water balance and the crop water requirement through the year (screen captures of release 1.1.0).}\label{fig:shot-real-balance}
\end{figure}

\begin{figure}[H]\centering
\includegraphics[width=0.62\textwidth,height=0.42\textheight,keepaspectratio]{\figdir/shots_jpg/real_tillage_flux.jpg}
\caption{SWAFL Desktop: pore-scale drying of the topsoil under conventional and no tillage: the evaporation flux and the suction at 1.25 and 3.75\,cm, and the distribution, velocity and content of the water on day 5.1 (screen capture of release 1.1.0).}\label{fig:shot-real-flux}
\end{figure}
